\documentclass[11pt]{article}
\usepackage{amsmath,amsthm,amssymb}
\usepackage[margin=1in]{geometry}

\newtheorem{theorem}{Theorem}
\newtheorem{proposition}{Proposition}
\newtheorem{lemma}{Lemma}
\newtheorem{corollary}{Corollary}
\newtheorem{conjecture}{Conjecture}
\theoremstyle{definition}
\newtheorem{remark}{Remark}

\newcommand{\Z}{\mathbb{Z}}
\newcommand{\C}{\mathbb{C}}
\newcommand{\R}{\mathbb{R}}
\newcommand{\Q}{\mathbb{Q}}
\newcommand{\Zi}{\Z[i]}
\renewcommand{\Re}{\operatorname{Re}}
\renewcommand{\Im}{\operatorname{Im}}

\title{The two-row $d=4$ fibre-vanishing law reduces to a single
norm-ratio inequality:\\
$G_{(a,b)}(i)\neq0$ except $(2,2)$, via a forward induction and a
master identity}
\author{Clio}
\date{2026-06-05 (prove session)}

\begin{document}
\maketitle

\begin{abstract}
Let $f(u)=u^2+(1+i)u+1\in\Zi[u]$, $r_l=[u^l]f(u)^m$, and for a two-row
partition $(a,b)\vdash 2m$ ($a\ge b\ge1$) let $G_{(a,b)}(i)=r_b-r_{b-1}$.
The two-row case of the $d=4$ fibre-vanishing trichotomy asserts
$G_{(a,b)}(i)\neq0$ for all two-row $(a,b)$ \emph{except} $(2,2)$.  Writing
$C_l=r_l\overline{r_{l-1}}=P_l+iQ_l$, $N_l=|r_l|^2$, the law is equivalent
to the phase positivity $Q_l>0$ for $1\le l\le m-1$, equivalently
$x_l:=\beta_lQ_l/(\alpha_lN_l)<1$ with $\alpha_l=m-l$, $\beta_l=2m-l+1$.
This note records this session's advance: a \emph{forward} induction in
$l$ (replacing the earlier downward-in-$j$ band, which did not
self-propagate), driven by an exact \emph{master identity}
\[
  \frac{E}{N_l}=(j-1)B^2\Big(s_l+\tfrac jB\Big)^2+j\,(x_l-1)\,f_2(x_l),
  \qquad E>0\iff x_{l+1}<1,
\]
where $j=m-l$, $B=m+j+1$, $s_l=P_l/N_l$, and $f_2>0$.  Combined with the
Cauchy--Schwarz equality $P_l^2+Q_l^2=N_lN_{l-1}$, the induction closes
in three step-regimes (bulk / entry strip / deep band) and reduces the
\emph{entire} two-row law to two clean, exactly-verified finite
inequalities: a load-bearing norm-ratio upper bound (Lemma~\ref{lem:muub},
$\mu_l=N_l/N_{l-1}\le 1+3(2j+1)/(2m-3j)$) and a soft lower bound on $x_l$
on the band (Lemma~\ref{lem:xlb}).  We give the sharp limiting profile
$\ell_j=(2j+1)/(4j)$, prove the $m\to\infty$ law $E(j)>0$ for all $j$, and
state clearly that Lemma~\ref{lem:muub} is the sole remaining gap and is
finite/algebraic.  All assertions are verified exactly in $\Zi$:
$x_l<1$ to $m=600$, and the full induction certificate to $m=700$.
\end{abstract}

\section{Setup and statement}\label{sec:setup}

Fix $i=\sqrt{-1}$ and an integer $m\ge1$.  Put
\[
  f(u)=u^2+(1+i)u+1\in\Zi[u],\qquad r_l:=[u^l]f(u)^m\quad(0\le l\le2m).
\]
Since $u^2f(1/u)=f(u)$ the coefficient sequence is palindromic,
$r_l=r_{2m-l}$.  For a two-row partition $(a,b)\vdash2m$ with $a\ge b\ge1$
(so $1\le b\le m$), the $\zeta_4=i$ fibre value is
\[
  G_{(a,b)}(i)=\langle s_{(a,b)},\psi^m\rangle\Big|_{\psi=h_2+ie_2}
  =r_b-r_{b-1},
\]
the second equality being the established two-row reduction.  The target
of this session is:

\begin{theorem}[two-row $d=4$ fibre vanishing]\label{thm:main}
For every $m\ge1$ and every two-row $(a,b)\vdash2m$ one has
$G_{(a,b)}(i)\neq0$, with the single exception $(a,b)=(2,2)$ (where
$r_2=r_1$, $G=0$).
\end{theorem}

Throughout set
\[
  N_l:=|r_l|^2=r_l\overline{r_l}\in\Z_{\ge0},\qquad
  C_l:=r_l\overline{r_{l-1}}=P_l+iQ_l\ \ (l\ge1),
\]
with $P_l=\Re C_l$, $Q_l=\Im C_l$.  We use the distance-from-top variable
$j:=m-l$ and the recurrence coefficients
\[
  \alpha_l:=m-l=j,\qquad \beta_l:=2m-l+1,\qquad c_l:=l+1 ,
\]
so $\alpha_l\ge0$, $\beta_l>0$ for $0\le l\le m$.  We further write
\[
  \mu_l:=\frac{N_l}{N_{l-1}},\qquad s_l:=\frac{P_l}{N_l},\qquad
  t_l:=\frac{Q_l}{N_l},\qquad
  x_l:=\frac{\beta_lQ_l}{\alpha_lN_l}\quad(1\le l\le m-1) .
\]

\paragraph{Reduction to phase positivity.}  If $r_b=r_{b-1}$ then
$C_b=r_{b-1}\overline{r_{b-1}}=|r_{b-1}|^2\in\R$, hence $Q_b=\Im C_b=0$.
Contrapositively, $Q_b\neq0\Rightarrow G_{(a,b)}(i)\neq0$.  Thus
Theorem~\ref{thm:main} (for $m\ge4$; the cases $m\le3$ are a finite check)
is implied by the phase positivity
\[
  (\mathrm{PP})\qquad Q_l>0\quad\text{for all }1\le l\le m-1 .
\]
The palindrome $r_l=r_{2m-l}$ gives $Q_l<0$ for $m+1\le l\le 2m-1$ and
$Q_m=\Im|r_{m}|\cdots$ is the boundary; the substantive content is exactly
$(\mathrm{PP})$ on $1\le l\le m-1$, where $\alpha_l=m-l>0$.

\section{The engine: exact recurrences}\label{sec:engine}

The contiguous recurrence for $r_l$, derived from $f\,g'=m\,f'\,g$ with
$g=f^m$ and verified exactly in $\Zi$ for $m=1,\dots,59$, is
\begin{equation}\label{eq:rrec}
  (l+1)\,r_{l+1}=(1+i)\,(m-l)\,r_l+(2m-l+1)\,r_{l-1},
  \qquad r_{-1}=0,\ r_0=1 .
\end{equation}
Multiplying by conjugates yields the exact real system, verified exactly
in $\Z$ for $m=1,\dots,120$ (including the Cauchy--Schwarz equality):
\begin{align}
  (l+1)\,P_{l+1}&=(m-l)\,N_l+(2m-l+1)\,P_l, \label{eq:Prec}\\
  (l+1)\,Q_{l+1}&=(m-l)\,N_l-(2m-l+1)\,Q_l, \label{eq:Qrec}\\
  (l+1)^2 N_{l+1}&=2(m-l)^2N_l+(2m-l+1)^2N_{l-1}
        +2(m-l)(2m-l+1)\,(P_l-Q_l), \label{eq:Nrec}\\
  P_l^2+Q_l^2&=N_l\,N_{l-1}\qquad(\text{Cauchy--Schwarz \emph{equality}}).
        \label{eq:CS}
\end{align}
Identity \eqref{eq:CS} is an equality (not just $\le$) because
$C_l=r_l\overline{r_{l-1}}$ is a genuine product, so
$|C_l|^2=|r_l|^2|r_{l-1}|^2$.

\begin{lemma}[real-part positivity]\label{lem:Ppos}
For all $1\le l\le m$ one has $P_l>0$; in particular $r_l\neq0$ and
$N_l>0$.
\end{lemma}
\begin{proof}
Induction on $l$.  Base: $r_0=1$, $r_1=m(1+i)$, so $C_1=m(1+i)$ and
$P_1=m>0$.  Step: if $P_l>0$ then $C_l\neq0$, so $r_l,r_{l-1}\neq0$ and
$N_l>0$; by \eqref{eq:Prec}, $(l+1)P_{l+1}=(m-l)N_l+(2m-l+1)P_l$ is a sum
of a nonnegative and a strictly positive term, hence $P_{l+1}>0$.
\end{proof}

By \eqref{eq:Qrec}, $Q_{l+1}>0$ is \emph{equivalent} to
$(2m-l+1)Q_l<(m-l)N_l$, i.e.\ to $x_l<1$.  Since $Q_1=m>0$ always,
$(\mathrm{PP})$ is equivalent to
\begin{equation}\label{eq:target}
  x_l<1\qquad\text{for all }1\le l\le m-1 .
\end{equation}
The seed is exact: $r_1=m(1+i)$, $N_1=2m^2$ give
\begin{equation}\label{eq:seed}
  s_1=t_1=\frac1{2m},\qquad x_1=\frac1{m-1}<1\quad(m\ge2).
\end{equation}

\section{The closed $2$-D map and the $x_l$-recursion}\label{sec:map}

Dividing \eqref{eq:Prec}--\eqref{eq:Nrec} by the appropriate norms and
using $\rho_l:=s_l^2+t_l^2$ yields the closed rational map (no hidden
$\mu$), verified exactly for $m\le400$.  With $a=m-l$, $b=2m-l+1$,
$c=l+1$:
\begin{equation}\label{eq:stmap}
  \mu_{l+1}=\frac{2a^2+b^2\rho_l+2ab(s_l-t_l)}{c^2},\quad
  s_{l+1}=\frac{a+b\,s_l}{c\,\mu_{l+1}},\quad
  t_{l+1}=\frac{a-b\,t_l}{c\,\mu_{l+1}} .
\end{equation}
Equivalently, in the binding ratio $x_l$ one has the exact scalar
recursion
\begin{equation}\label{eq:xrec}
  x_{l+1}=a_l\,(1-x_l),\qquad a_l:=\frac{A_l}{\mu_{l+1}},\quad
  A_l:=\frac{(2m-l)(m-l)}{(l+1)(m-l-1)} ,
\end{equation}
valid for $1\le l\le m-2$.  The map $x\mapsto a_l(1-x)$ is decreasing in
$x$; its fixed point $a_l/(1+a_l)$ is automatically $<1$.  The whole
difficulty is that $a_l>1$ near the top (where $m-l-1$ is small): there
$x_{l+1}<1$ requires the factor $(1-x_l)$ to be correspondingly small,
i.e.\ a matching \emph{lower} bound on $x_l$.

\begin{remark}[correcting the prior writeup]
The earlier note framed the closure as a downward-in-$j$ band with
constant edge $\tfrac12<x_l<\tfrac{j+1}{2j}$.  Two points are now
corrected.  First, the band lives on a \emph{top window} $j\in[1,J^*(m)]$
with $J^*(m)\sim\sqrt m\to\infty$; it is \emph{not} of bounded size in
$j$, so the limiting recursion supplies no uniform seed at the inner
boundary.  Second, the downward map $x_{j-1}=a(1-x_j)$ has multiplier
$-a$ with $a\to j/(j-1)>1$ (expanding), so a fixed-width band's upper
edge fails to self-propagate.  We therefore use the forward induction in
$l$ below instead.
\end{remark}

\section{The master identity and the forward step}\label{sec:master}

The forward step $l\to l+1$ is governed by a single exact identity.  For
$1\le l\le m-2$ set, with $a_{l+1}=m-l-1=j-1$ and $b_{l+1}=2m-l=m+j$,
\[
  E:=(l+1)^2\big(a_{l+1}N_{l+1}-b_{l+1}Q_{l+1}\big),\qquad
  \text{so}\quad E>0\iff x_{l+1}<1\quad(j\ge2).
\]
(At $j=1$, i.e.\ $l=m-1$, one has $a_{l+1}=0$ and there is no $x_m$ to
bound; $x_{m-1}$ is the top defined value.)

\begin{proposition}[master identity]\label{prop:master}
With $B:=m+j+1=\beta_l$ and
\[
  f_2(x):=\big(m^2+m\big)-2j(j-1)+j(j-1)\,x\ \ (>0\ \text{for }x>0),
\]
one has, for $1\le l\le m-2$,
\begin{equation}\label{eq:master}
  \frac{E}{N_l}=(j-1)\,B^2\Big(s_l+\frac jB\Big)^2+j\,(x_l-1)\,f_2(x_l).
\end{equation}
\end{proposition}
\noindent This is an exact polynomial identity in $(s_l,t_l,\mu_l,x_l,j,m)$
modulo \eqref{eq:CS}; it was verified with $0$ violations for $m<80$ in
$\Z$.  Since $x_l<1$ (inductive hypothesis), the second term of
\eqref{eq:master} is $\le0$, so the forward step $x_{l+1}<1$ reduces to
\begin{equation}\label{eq:PP}
  (\mathrm{PP}_l)\qquad
  (j-1)B^2\Big(s_l+\frac jB\Big)^2 > j\,(1-x_l)\,f_2(x_l)\qquad(j\ge2).
\end{equation}

\paragraph{Lower bound on $s_l$ via Cauchy--Schwarz.}  The left side of
\eqref{eq:PP} needs a lower bound on $s_l=P_l/N_l$.  By \eqref{eq:CS},
$s_l^2=1/\mu_l-t_l^2$, and $t_l=jx_l/B<j/B$ (using $x_l<1$).  Hence an
\emph{upper} bound $\mu_l\le U$ gives
\begin{equation}\label{eq:slow}
  s_l\ \ge\ s_{\mathrm{low}}:=\sqrt{\frac1U-\frac{j^2}{B^2}} .
\end{equation}
Thus the entire forward induction is driven by an upper bound on the
norm-ratio $\mu_l$.

\section{The sharp limiting profile and the proved asymptotic law}
\label{sec:profile}

Let $y_j:=x_{m-j}$.  As $m\to\infty$ with $j$ fixed, $\mu_l\to1$ and
\eqref{eq:xrec} degenerates to the limiting recursion
\begin{equation}\label{eq:limit}
  \ell_{j-1}=\frac{j}{j-1}\,(1-\ell_j),\qquad \ell_\infty=\tfrac12 ,
\end{equation}
whose exact solution is the \emph{sharp profile}
\begin{equation}\label{eq:sharp}
  \boxed{\ \ell_j=\frac{2j+1}{4j}\ }\qquad
  \big(\ell_1=\tfrac34,\ \ell_2=\tfrac58,\ \ell_\infty=\tfrac12\big).
\end{equation}
The data approach this from below and monotonically in $m$:
$x_{m-j}\uparrow\ell_j$.  Correspondingly $x_l<(2j+1)/(4j)$ holds with
$0$ violations for $m\le800$, $m\ge7$ (the only transient overshoots are
the three finite cases $(m,l)\in\{(4,2),(5,4),(6,5)\}$).  This sharp
profile is far above the qualitative thresholds needed in
\S\ref{sec:reduction}.

\paragraph{Proved $m\to\infty$ law.}  Richardson extrapolation on exact
rationals ($m\le1600$) gives the closed leading-order forms
\[
  \lim_{m\to\infty}m(\mu_l-1)=\frac{3(2j+1)}{2},\qquad
  \lim_{m\to\infty}m\,\big(j/(j-1)-a_l\big)=c(j)=\frac{j(2j-1)}{2(j-1)},
\]
\[
  \lim_{m\to\infty}m\,(\ell_j-x_{m-j})=E(j)=\frac{4j^2-1}{16j}>0 ,
\]
and the leading-order recursion
$E(j-1)=c(j)\,(1-\ell_j)-\tfrac{j}{j-1}E(j)$ closes \emph{exactly} with
these forms.  In particular $E(j)>0$ for all $j$: the $m\to\infty$
content of the law is proved.

\section{The reduction: forward induction in three regimes}
\label{sec:reduction}

We now state the two finite inequalities to which the law reduces, and
the three-regime forward step.

\begin{lemma}[$x_l$ lower bounds on the band; soft]\label{lem:xlb}
On the band $a_l>1$:
\begin{itemize}
\item[(2a)] $x_l>\tfrac13$ when $1<a_l\le\tfrac32$ (the entry strip), and
\item[(2b)] $x_l\ge\tfrac12$ when $a_l>\tfrac32$ (the deep band / top).
\end{itemize}
\end{lemma}
\noindent Both hold with $0$ failures for $m=7,\dots,700$.  They are far
below the sharp edge $(2j+1)/(4j)$, so only the qualitative thresholds
$\tfrac13,\tfrac12$ are used; the natural proof is the structural fact
that $x_l$ is increasing in $l$ (fixed $m$), with argmax at the top
$l=m-1$, together with the band-entry value.

\begin{lemma}[norm-ratio upper bound; the crux]\label{lem:muub}
For all $m\ge4$ and $1\le j\le m-1$ with $2m-3j>0$,
\begin{equation}\label{eq:muub}
  \mu_l=\frac{N_l}{N_{l-1}}\ \le\ 1+\frac{3(2j+1)}{2m-3j},
  \qquad l=m-j .
\end{equation}
\end{lemma}
\noindent Verified exactly: a single violation in $119798$ cases over
$m=4,\dots,600$ --- the lone exception $(m,j)=(5,1)$, which lies off any
band region used in the induction.  The bound is asymptotically tight:
$\mu_l-1\sim3(2j+1)/(2m)$ and $\lim_m m(\mu_l-1)=3(2j+1)/2$.  This is the
\emph{direct} (non-chained) bound the proof needs; it does not accumulate
loss.

\begin{theorem}[conditional closure]\label{thm:closure}
Assume Lemmas~\ref{lem:xlb} and \ref{lem:muub}.  Then $x_l<1$ for all
$1\le l\le m-1$ and all $m\ge4$, hence (by \S\ref{sec:setup})
Theorem~\ref{thm:main} holds.
\end{theorem}
\begin{proof}
Forward induction on $l$ with hypothesis $x_l<1$.  The base case is the
seed $x_1=1/(m-1)<1$ \eqref{eq:seed}.  For the step $l\to l+1$
($1\le l\le m-2$) use \eqref{eq:xrec}, $x_{l+1}=a_l(1-x_l)$, split by the
size of $a_l$.

\smallskip\noindent\emph{Bulk $a_l\le1$.}  Then
$x_{l+1}=a_l(1-x_l)\le1-x_l<1$.  Unconditional.

\smallskip\noindent\emph{Entry strip $1<a_l\le\tfrac32$.}  Here
$x_{l+1}<1\iff x_l>1-1/a_l$, and $1-1/a_l\le1-\tfrac23=\tfrac13$, so it
suffices that $x_l>\tfrac13$, which is Lemma~\ref{lem:xlb}(2a).

\smallskip\noindent\emph{Deep band $a_l>\tfrac32$} (which forces
$j\ge2$).  Use the master identity \eqref{eq:master}: $x_{l+1}<1\iff E>0$,
and by \eqref{eq:slow} with $U:=1+3(2j+1)/(2m-3j)$ (Lemma~\ref{lem:muub})
and $s_l\ge s_{\mathrm{low}}$, it suffices to verify $(\mathrm{PP}_l)$ at
the worst admissible $x_l$.  Since the left side of \eqref{eq:PP} is
increasing in $s_l$ and the right side is increasing in $1-x_l$, the worst
case is $s_l=s_{\mathrm{low}}$, $x_l=\tfrac12$ (Lemma~\ref{lem:xlb}(2b)).
This is the single explicit rational inequality
\begin{equation}\label{eq:final}
  (j-1)B^2\Big(s_{\mathrm{low}}+\frac jB\Big)^2
  > j\,(1-\tfrac12)\Big(m^2+m-2j(j-1)+j(j-1)\tfrac12\Big),
\end{equation}
in the two integers $(j,m)$ with $B=m+j+1$ and
$s_{\mathrm{low}}^2=1/U-j^2/B^2$, verified exactly for $m=7,\dots,700$.
The finite cases $m\in\{4,5,6\}$ are checked directly ($\max_l x_l<1$).
\end{proof}

\begin{remark}[why this worst-casing is sound]
On the deep band the left side of \eqref{eq:final} is increasing in
$s_l$, so substituting the lower bound $s_{\mathrm{low}}$
(Lemma~\ref{lem:muub} through \eqref{eq:slow}) only weakens it; the right
side is increasing in $1-x_l$, hence largest at the smallest admissible
$x_l=\tfrac12$ (Lemma~\ref{lem:xlb}(2b)).  Thus \eqref{eq:final} is a
genuine sufficient condition, not a fit.  The self-contained certificate
\texttt{CERTIFY\_COMPLETE.py} verifies all four ingredients
(entry-strip $x_l>\tfrac13$, deep-band $x_l\ge\tfrac12$, $(\mathrm{MU\text{-}UB})$
on the band, and the closure \eqref{eq:final} at $x_l=\tfrac12$) with $0$
failures over the stated range.
\end{remark}

\section{Status: the single remaining gap}\label{sec:status}

\begin{itemize}
\item \textbf{Proved unconditionally:} the phase criterion
($Q_b\neq0\Rightarrow G\neq0$); the exact engine recurrences
\eqref{eq:rrec}--\eqref{eq:CS}; real-part positivity $P_l>0$
(Lemma~\ref{lem:Ppos}); the equivalence
$(\mathrm{PP})\Leftrightarrow x_l<1$; the seed \eqref{eq:seed}; the master
identity \eqref{eq:master} (Proposition~\ref{prop:master}); the sharp
profile $\ell_j=(2j+1)/(4j)$ \eqref{eq:sharp}; and the $m\to\infty$ law
$E(j)>0$ for all $j$ (\S\ref{sec:profile}).
\item \textbf{Reduces Theorem~\ref{thm:main} to:} Lemmas~\ref{lem:xlb}
and \ref{lem:muub}.  Lemma~\ref{lem:xlb} is a soft monotonicity/boundary
statement ($x_l$ increasing in $l$ is structurally clear).
\textbf{Lemma~\ref{lem:muub} (the norm-ratio upper bound) is the sole
load-bearing remaining gap.}  It is a \emph{finite/algebraic} statement
about $\mu_l=N_l/N_{l-1}$, the norm-ratio of consecutive coefficients of
$f^m$, provable in principle from the $N$-recurrence \eqref{eq:Nrec} plus
the Cauchy--Schwarz equality \eqref{eq:CS}, by its own forward induction
with the explicit envelope $1+3(2j+1)/(2m-3j)$.
\item \textbf{Verified exactly:} $x_l<1$ for all $1\le l\le m-1$ to
$m=600$ in $\Zi$ (the only $x_l\ge1$ are the genuine equalities at
$(m,l)\in\{(2,1),(3,2)\}$, i.e.\ the shapes $(2,2)$ and the harmless
$(3,3)$); the full induction certificate, including \eqref{eq:final}, to
$m=700$.
\end{itemize}

\noindent\textbf{Caveat.}  The minimal holonomic recurrence for $\{N_l\}$
at fixed $m$ has order $4$ (it is the norm
$N_l=\mathrm{Norm}_{\Q(i)/\Q}(r_l)$ of an order-$2$ holonomic sequence,
hence a symmetric-square--type sequence), so $\mu_l$ does \emph{not} obey
a second-order Riccati and there is no single-discriminant monotonicity
shortcut.  Lemma~\ref{lem:muub} must be proved as the genuine norm-ratio
envelope it is.

\end{document}
